\documentclass[a4paper,11pt]{article}
\usepackage{exam-style-341}
\examinehead{341 农业综合知识三 · 网络技术分章练习}{第05章 参考答案与解析}
\begin{document}

\unitanswerscover{341 农业综合知识三 · 网络技术分章练习}{第05章参考答案}{局域网与广域网：逐题答案与解析；争用期、最短帧长与退避时间计算题给出完整步骤与单位。}

\answerpart{一、单项选择题（10 题，每题 3 分，共 30 分）}

\q A。\begin{analysis}IEEE 802.3 是以太网标准；802.11 是无线局域网标准；802.15 是无线个人区域网标准；802.16 是城域网标准（教材 3.3.1，印刷页 84—88）。\end{analysis}

\q A。\begin{analysis}CSMA/CD 的要点是多点接入、载波监听和碰撞检测，可概括为“先听后发、边发边听、冲突停止、随机重发”（教材 3.3.2，印刷页 85—92）。B 属于网络层／运输层，C 属于物理层，D 属于数据链路层的帧处理。\end{analysis}

\q B。\begin{analysis}10 Mb/s 以太网的争用期取 $51.2\ \mu\text{s}$，即 512 比特时间，它是端到端往返传播时延的 2 倍（教材 3.3.2，印刷页 90—91）。\end{analysis}

\q B。\begin{analysis}10 Mb/s 以太网的最短帧长是 64 B，即 512 bit；发送 512 bit 所需时间正好是争用期 51.2 $\mu$s（教材 3.3.2，印刷页 91）。\end{analysis}

\q B。\begin{analysis}截断二进制指数退避算法：第 $i$ 次冲突后取 $k=\min(i,10)$，从整数集合 $\{0,1,\dots,2^k-1\}$ 中随机取一个数 $r$，退避时间为 $r$ 倍争用期（教材 3.3.2，印刷页 90）。\end{analysis}

\q D。\begin{analysis}MAC 地址（硬件地址）长度为 48 bit，即 6 B，固化在网络适配器中（教材 3.3.5）。\end{analysis}

\q A。\begin{analysis}交换机每收到一个帧，就记录帧的源 MAC 地址与接收端口，据此建立／更新 MAC 地址表；转发时按目的 MAC 地址查表（教材 3.4.2）。B 目的地址用于转发查找，不是学习依据。\end{analysis}

\q B。\begin{analysis}802.11 无线局域网的 MAC 层使用 CSMA/CA（载波监听多点接入／碰撞避免）协议（教材 9.1.3，印刷页 412）。\end{analysis}

\q A。\begin{analysis}100BASE-T 的含义：100 表示数据率 100 Mb/s，BASE 表示基带传输，T 表示双绞线（Twisted pair）（教材 3.5.1，印刷页 106）。\end{analysis}

\q A。\begin{analysis}PPP 是广域网中广泛使用的点对点数据链路层协议，用于路由器之间、用户与 ISP 之间的链路（教材 3.2，印刷页 78—84）。HTTP、UDP 分别属于应用层和运输层，CSMA/CA 用于无线局域网。\end{analysis}

\answerpart{二、多项选择题（3 题，每题 4 分，共 12 分）}

\q ABC。\begin{analysis}以太网 MAC 帧包括目的 MAC 地址、源 MAC 地址、类型字段、数据字段和 FCS；路由表是路由器转发分组时使用的，不是以太网帧的字段（教材 3.3.5）。\end{analysis}

\q ABC。\begin{analysis}A、B、C 都是 CSMA/CD 的正确描述；D 错误：全双工方式下发送和接收使用不同的信道，不会发生冲突，因此不使用 CSMA/CD（教材 3.3.2、3.5）。\end{analysis}

\q ABC。\begin{analysis}A、B 正确：802.11 无线局域网使用 CSMA/CA；C 正确：可用 RTS/CTS 机制缓解隐藏站问题；D 错误：无线局域网不能像有线以太网那样边发送边检测冲突，因此不使用 CSMA/CD（教材 9.1.3，印刷页 412—418）。\end{analysis}

\answerpart{三、判断题（6 题，每题 2 分，共 12 分）}

\q ×。改正：10 Mb/s 以太网的争用期是 51.2 $\mu$s（即 512 比特时间），不是 512 $\mu$s。\begin{analysis}争用期取端到端往返传播时延的 2 倍，10 Mb/s 以太网规定为 51.2 $\mu$s，对应发送 512 bit（64 B）的时间（教材 3.3.2，印刷页 90—91）。\end{analysis}

\q √。\begin{analysis}以太网最短帧长为 64 B；发送数据时若帧长不足 64 B 需填充。长度小于 64 B 的帧都是因冲突而异常中止的无效帧（教材 3.3.2，印刷页 91）。\end{analysis}

\q ×。改正：交换机自学习建立 MAC 地址表时，依据的是帧的源 MAC 地址和接收端口。\begin{analysis}交换机通过记录“源 MAC 地址—入端口”的对应关系建立地址表；按目的 MAC 地址查表转发。本题把学习依据与转发依据混淆了（教材 3.4.2）。\end{analysis}

\q √。\begin{analysis}VLAN 通过逻辑划分把一个物理局域网分成多个广播域，可减少广播风暴、提高安全性和管理灵活性（教材 3.4.3，印刷页 104）。\end{analysis}

\q ×。改正：全双工以太网不再使用 CSMA/CD，因为收发使用不同信道，不会发生冲突。\begin{analysis}CSMA/CD 只用于半双工共享信道；全双工方式下发送和接收可同时进行，无需进行冲突检测（教材 3.5）。\end{analysis}

\q √。\begin{analysis}无线环境中信号强度动态范围大、存在隐藏站问题，发送时无法有效检测是否发生冲突，因此 802.11 采用 CSMA/CA，通过预约信道、确认和退避来避免冲突（教材 9.1.3，印刷页 412—416）。\end{analysis}

\answerpart{四、计算题（2 题，每题 5 分，共 10 分）}

\q 参考答案：
（1）最大单程传播时延 $t=2000\ \text{m}\div(2\times10^{8}\ \text{m/s})=1\times10^{-5}\ \text{s}=10\ \mu\text{s}$（1.5 分）。
（2）最小争用期取往返传播时延：$2t=2\times10\ \mu\text{s}=20\ \mu\text{s}$（1.5 分）。
（3）最短帧长 $=2t\times$ 数据率 $=20\times10^{-6}\ \text{s}\times100\times10^{6}\ \text{bit/s}=2000\ \text{bit}=250\ \text{B}$（2 分）。
\begin{analysis}最短帧长的作用是保证发送方在帧发送完毕之前就能检测到冲突：帧的发送时间必须不小于争用期。10 Mb/s 以太网在同样 2 km 跨度下只需 $20\ \mu\text{s}\times10\ \text{Mb/s}=200\ \text{bit}=25\ \text{B}<64\ \text{B}$，所以标准取 64 B 更保守，可覆盖约 5 km 的跨度；而 100 Mb/s 以太网若仍要覆盖 2 km，就需要 250 B 的最短帧长，这正是高速以太网要限制网络直径或改用全双工的原因。\end{analysis}

\q 参考答案：
（1）第 3 次冲突后 $k=\min(3,10)=3$（1 分）。
（2）随机数 $r$ 的取值范围为 $\{0,1,2,3,4,5,6,7\}$，即 $0\leq r\leq2^3-1=7$（2 分）。
（3）退避时间 $=r\times$ 争用期 $=r\times51.2\ \mu\text{s}$，故可能范围为 $0\sim7\times51.2=358.4\ \mu\text{s}$，可取 $0,51.2,102.4,\dots,358.4\ \mu\text{s}$（2 分）。
\begin{analysis}截断二进制指数退避使冲突次数越多、可选的退避范围越大，从而降低再次冲突的概率；$k$ 最大取 10，即退避范围最大为 $0\sim1023$ 倍争用期。本题注意“第 3 次冲突”对应 $k=3$ 而不是 $k=10$。\end{analysis}

\answerpart{五、简答题（2 题，每题 8 分，共 16 分）}

\q 参考答案要点：
（1）工作过程（5 分）：CSMA/CD 即载波监听多点接入／碰撞检测。（1） 多点接入：多个站点连接在同一条总线上；（2） 载波监听：发送前先监听信道，若忙则等待，若闲则发送（先听后发）；（3） 碰撞检测：发送过程中继续监听，若检测到冲突则立即停止发送，并发送人为干扰信号（边发边听、冲突停止）；（4） 随机重发：按截断二进制指数退避算法等待一段随机时间后重新尝试发送（随机重发）（教材 3.3.2，印刷页 85—92）。
（2）争用期与最短帧长的关系（3 分）：争用期是端到端往返传播时延的 2 倍，是发送方能够检测到冲突的最长时间；为保证发送方在帧发送完毕之前能检测到冲突，帧的发送时间必须不小于争用期，即 最短帧长 $=$ 争用期 $\times$ 数据率。10 Mb/s 以太网争用期 51.2 $\mu$s，最短帧长 $=51.2\times10^{-6}\times10\times10^{6}=512\ \text{bit}=64\ \text{B}$。小于 64 B 的帧被视为冲突碎片而丢弃。

\q 参考答案要点：
（1）地址学习（3 分）：交换机收到帧后，把帧的源 MAC 地址与接收端口号写入 MAC 地址表（若表中已有则刷新），并设置老化时间（教材 3.4.2）。
（2）转发与过滤（3 分）：交换机按帧的目的 MAC 地址查表：若查到对应端口且不是入端口，则从该端口转发；若查到端口就是入端口，则丢弃（过滤）；若查不到，则向除入端口外的所有端口泛洪。这样可隔离碰撞域、提高网络性能。
（3）VLAN 的作用（2 分）：VLAN 把一个物理局域网逻辑上划分为多个广播域，可限制广播风暴、提高网络安全性、便于按业务或部门灵活管理；不同 VLAN 之间通信需要路由器或三层交换机（教材 3.4.3，印刷页 104）。

\answerpart{六、综合题（1 题，共 20 分）}

\q 参考答案要点：
（1）（4 分）交换机自学习：收到帧时记录源 MAC 地址与入端口，建立“MAC 地址—端口”映射并定期老化；收到目的 MAC 地址未知的帧时，向除入端口外的所有端口泛洪转发，待目的站回复后即可学习到其端口。
（2）（4 分）VLAN 的作用：把一个物理局域网划分为多个广播域，减少广播风暴、提高安全性和管理灵活性。不同 VLAN 之间不能直接二层通信，需要路由器或三层交换机进行三层转发。
（3）（4 分）温室无线终端采用 IEEE 802.11（Wi-Fi）标准，MAC 层使用 CSMA/CA。无线环境不能使用 CSMA/CD，因为无线站点发送时无法同时可靠接收并检测冲突（信号强度动态范围大、存在隐藏站和暴露站问题），故采用碰撞避免机制：先监听、随机退避，必要时用 RTS/CTS 预约信道，并用 ACK 确认（教材 9.1.3）。
（4）（4 分）发送时延 $=1000\ \text{B}\times8\div(100\times10^{6}\ \text{bit/s})=8000\ \text{bit}\div(100\times10^{6}\ \text{bit/s})=8\times10^{-5}\ \text{s}=80\ \mu\text{s}$；
传播时延 $=10\ \text{km}\div(2\times10^{8}\ \text{m/s})=10^{4}\ \text{m}\div(2\times10^{8}\ \text{m/s})=5\times10^{-5}\ \text{s}=50\ \mu\text{s}$。
高速以太网采用全双工方式：收发使用不同的双绞线对或光纤，发送时不会发生冲突，因而不再使用 CSMA/CD；同时可避免“争用期与最短帧长”的矛盾，提高链路利用率和有效带宽。
（5）（4 分）局域网覆盖范围小、通常为一个单位所有，核心技术是以太网和无线局域网，典型协议为 IEEE 802.3、802.11；广域网覆盖范围大，通常借用公用通信网，核心技术是点对点链路和分组交换，典型协议有 PPP、HDLC、帧中继、ATM 等。园区网范围小、要求高带宽低时延、设备成本低且便于管理，因此选择交换式以太网（全双工）技术；广域网则需借助电信运营商的链路和相应协议实现远距离互连。

\end{document}
